\section{Results}\label{sec:results}
\subsection{Study 1 in brief: execution-aware stopping on the nominal model}\label{sec:study1}
Study~1 made only the stopping rule execution-aware and kept the nominal planning model. In \val{JNRunsMain} tier B
runs, AnytimeSwarm-R needed \val{JMsRedAS}\% less computation per call than AnytimeSwarm, but its advantage over ALNS
with 40\,ms was \val{JDFixedForty}\,percentage points (pp) with a Holm-adjusted $p$ of \val{JPFixedForty}, and its
advantages over AnytimeSwarm and the stagnation rule were not significant either. Periodic replanning raised the
utility of every method by \val{JTrigMin} to \val{JTrigMax}\,pp. Supplementary Sections~S3 and~S4 report these results
in full, together with an exact-execution benchmark, tier C flights and a scalability study with 40 UAVs. Study~2 was
designed after these results and moves the error model from the decision to stop into the planning model.

\subsection{Study 2: the chance-constrained model against baselines}\label{sec:study2}
Study~2 comprises \val{KNRuns} tier B runs, that is, \val{KNMissions} held-out missions in \val{KNConfigs}
configurations with \val{KNMethods} methods (Table~\ref{tab:s2main}). The proposed replanner
reached the highest utility, \val{KUProp}\%. Its advantage over the six baselines of H1 ranged from \val{KDmBuf}\,pp
over the constant buffer to \val{KDmForty}\,pp over ALNS with 40\,ms, and every Holm-adjusted $p$-value was below
0.001, so H1 holds for all six comparisons. In relative terms, it served \val{KURelBuf}\% more reward than the constant
buffer and \val{KURelOnl}\% more than online performance prediction. Against the conventional planner, the advantage
was \val{KDConv}\,pp, or \val{KURelConv}\%, so H2 holds as well.

The advantage comes from fewer late tasks. The proposed replanner served \val{KLateProp} tasks per mission after their
deadline, against \val{KLateForty} for ALNS with 40\,ms, that is, \val{KLateRedForty}\% fewer, and \val{KLateBuf} for
the constant buffer. It lost no task to a late plan, whereas ALNS with 40\,ms lost \val{KLostForty} and the full budget
of 250\,ms \val{KLostFix} tasks per mission. On the nominal model, the stopping rules hardly differed. ALNS with
40\,ms, the size-scaled budget, the performance profile, online prediction, AnytimeSwarm and AnytimeSwarm-R lay between
\val{KUForty}\% and \val{KURisk}\%, although their computation ranged from \val{KMsOnl} to \val{KMsAS}\,ms per call
(Supplementary Fig.~S4).

\inputtable{tab_s2_main}


The ablations confirm that the planning model carries the gain (Table~\ref{tab:s2main}, bottom). With a fixed budget
of 40\,ms instead of the proposed stopping rule, utility changed by \val{KDAblStop}\,pp. The 90\% interval of this
difference, \val{KTostCI}\,pp, lies inside the margin of $\pm0.5$\,pp, so the stopping rule is equivalent to the fixed
budget on the chance-constrained model (H3). It still guards against long budgets on slow planners. At $\kappa_B=200$,
the full budget lost \val{KLostFixKTwoHundred} tasks per mission to late plans and served \val{KDFixKTwoHundred}\,pp less
reward than the proposed replanner. With the event trigger instead of the periodic one, the proposed replanner lost
\val{KDAblTrig}\,pp, far less than the \val{JTrigMin} to \val{JTrigMax}\,pp that periodic replanning added on the nominal
model of study~1. With $z=0$, the difference was \val{KDAblZ}\,pp and not significant.

The gain was not uniform. Per configuration, the proposed replanner was better than the constant buffer in \val{KWinBuf}
of the \val{KNConfigs} configurations and worse in \val{KLossBuf}. All losses occurred under high disturbances with
halved windows, where it served \val{KDBufWHalfHigh}\,pp relative to the buffer and \val{KDFortyWHalfHigh}\,pp relative
to ALNS with 40\,ms. With short windows and large errors, the effective deadline of many tasks falls before the earliest
possible arrival, so the planner gives up tasks that a less cautious plan still serves in time. This weakness motivated
the fallback of Section~\ref{sec:fallback}.

\subsection{Study 3: nominal fallback for given-up tasks}\label{sec:study3}
Study~3 comprises \val{LNRuns} tier B runs on \val{LNMissions} new held-out missions in \val{LNConfigs} configurations
with five methods (Table~\ref{tab:s3}, Fig.~\ref{fig:s3}). It first replicates study~2. The proposed replanner of
study~2 served \val{LDPropRepBuf}\,pp more reward than the constant buffer over all configurations, and it again lost to
the buffer under high disturbances with halved windows, by \val{LDPropBufHighWHalf}\,pp. The weakness is therefore
systematic and not a property of the seeds of study~2.

The fallback variant reached a utility of \val{LUFB}\%, the highest of all methods, and all four primary comparisons
were significant (Holm-adjusted $p<0.001$). It served \val{LDPropAll}\,pp more reward than the proposed replanner of
study~2 (H5), that is, \val{LURelProp}\% more. Under high disturbances with halved windows, it served
\val{LDBufHighWHalf}\,pp more than the constant buffer (H6). Over all configurations, it served \val{LDFortyAll}\,pp more
than ALNS with 40\,ms (H7, \val{LURelForty}\%) and \val{LDBufAll}\,pp more than the constant buffer (H8,
\val{LURelBuf}\%). Against the conventional planner, the advantage was \val{LDConvAll}\,pp, or \val{LURelConv}\%. The gain
over the replanner of study~2 was concentrated where the chance constraint gives up tasks, with \val{LDPropHighWHalf}\,pp
under high disturbances with halved windows. With the original windows, the two variants were close
(\val{LDPropMidWOne}\,pp under mid and \val{LDPropHighWOne}\,pp under high disturbances).

\inputtable{tab_s3}

\begin{figure}[t]
\centering
\fig{fig_s3}{\linewidth}
\caption{Study~3, tier B, held-out seeds 7000--7029: utility of the replanner with nominal fallback minus each other
method, per regime (mean over $\kappa_B$, \val{LNMissions} missions, 95\% paired bootstrap intervals).}\label{fig:s3}
\end{figure}

The fallback has a price in punctuality. A fallback task is planned against its nominal deadline and is often served
late. The fallback variant served \val{LLateFB} tasks late per mission, \val{LLateIncProp}\% more than the
\val{LLateProp} of the replanner of study~2, although still \val{LLateRedForty}\% fewer than ALNS with 40\,ms. It also
lost \val{LLostFB} tasks per mission to late plans, against none for the replanner of study~2. A late service yields no
reward, so these attempts cost only flight time. Where a late arrival has a cost of its own, the guarantee of
Proposition~\ref{prop:cc} without fallback remains the safer choice.

\subsection{Software-in-the-loop flights}\label{sec:s2-tierc}
All \val{KCRuns} held-out flights of study~2 were completed (Table~\ref{tab:s2tierc}). As specified in the protocol,
\val{KCInvalid} runs with a real-time factor outside the validity range were repeated once. Because the laptop was also
used for other work, many repeats ran slow as well, and two sensitivity analyses that were not pre-registered, with
the original runs and with valid pairs only, led to the same conclusions (Supplementary Section~S5).

The proposed replanner reached \val{KCUProp}\%, against \val{KCUForty}\% for ALNS with 40\,ms and \val{KCUConv}\% for
the conventional planner. The paired differences were \val{KCDForty}\,pp and \val{KCDConv}\,pp, with Holm-adjusted
$p$-values of \val{KCPForty} and \val{KCPConv}, so H4 is not supported. The mechanism of tier B reappeared in the counts. The proposed
replanner served \val{KCLateProp} tasks late per mission, against \val{KCLateForty} for ALNS with 40\,ms, which is
\val{KCLateRedForty}\% fewer, and it skipped no task at departure, whereas ALNS with 40\,ms skipped \val{KCSkipForty} and
the conventional planner \val{KCSkipConv} tasks per mission.

\inputtable{tab_s2_tierc}

The same missions were also run in tier B with the rescaled parameters. Tier B predicted the utility of a flight
closely ($r=$ \val{KBCr} over \val{KBCn} runs), but even tier B did not resolve the differences between the methods on
these \val{KCN} missions (\val{KBgzDForty}\,pp to ALNS with 40\,ms). With a per-mission standard deviation of
\val{KCSdForty}\,pp in tier C, confirming a difference of \val{KCEffForty}\,pp with a power of 80\% would require about
\val{KCNPowerForty} missions per method. Tier C thus shows that the replanner runs on autopilots with the same
computation and fewer late tasks, but it cannot confirm the size of the gain measured in tier B.
